\chapter{The PLATO Mission}
\label{chap:plato_mission}

\Cref{chap:kepler_tess} showed that the current transit catalogue does
not contain the planets the \HZIED test needs.  This chapter describes
the survey that is about to change that: ESA's \PLATO mission
\parencite{Rauer2016,ESA2017PLATO}, its target-sample structure
\parencite{Nascimbeni2022}, and the \PLATO Input Catalogue that feeds
the forecast in \Cref{chap:results}.


% Sections 2.6.1–2.6.4 merged into one section for condensed version
\section{PLATO mission and target samples}
\label{sec:plato_mission}
\phantomsection\label{sec:plato_target_samples}
\phantomsection\label{sec:pic_integration}
\phantomsection\label{sec:plato_samples}

\emph{Scaffold: mission concept.}
\PLATO as ESA's M3 medium-class mission and its role as the successor
to \Kepler{} and \TESS{}.  24 normal cameras and 2 fast cameras,
photometric precision $34\,\mathrm{ppm}\,\mathrm{h}^{-1/2}$ on the
brightest targets, structure of the long-duration observation phase
and two step-and-stare fields.

% Labels from former §Target samples P1-P5, §PLATO Input Catalogue integration,
% §PLATO Target Samples kept as phantoms above.

% Former \subsection{The PLATO Input Catalog and sample structure} merged:
\phantomsection\label{sec:pic}

The \PLATO Input Catalog \parencite{Montalto2021} is the official
pre-selected target list from which \PLATO will download light curves
during the long-duration observation phase.  The version used here
(PIC~v2.2.0.1) covers the LOPS2 southern field, whose coordinates were
identified by \textcite{Nascimbeni2022} as the optimal choice for the
first LOP field.

Stars in the PIC are grouped into four photometric priority samples,
formally defined in the \PLATO Science Requirements Document and
described in \textcite{Nascimbeni2022}. The numbering skips P3, which
was eliminated during the mission design phase while the remaining labels
were kept unchanged \parencite{Nascimbeni2022}.  The samples are defined
by spectral type, apparent magnitude and a minimum
photometric noise requirement for P1 and P2
(\Cref{tab:plato_samples}).

\begin{table}[H]
    \centering
    \caption[PLATO Input Catalog stellar samples]{%
        \PLATO stellar samples in the LOPS2 field (PIC~v2.2.0.1).
        Sample definitions from \textcite{Nascimbeni2022}; star counts
        from the processed PIC parquets used in this work; $\sigma_R$
        are per-sample pipeline defaults \parencite{ESA2017PLATO}.
        P3 was eliminated during mission design; the numbering is
        unchanged.}
    \label{tab:plato_samples}
    \begin{tabular}{lllrrl}
        \toprule
        Sample & Types & Limit & $N_\star$ (LOPS2) & $\sigma_R$ & Role \\
        \midrule
        P2 & F5--K7 dwarfs/subgiants & $V < 8.5$ & $605$     & $0.8\%$ & not analysed \\
        P1 & F5--K7 dwarfs/subgiants & $V < 11$  & $9\,665$  & $1\%$   & not analysed \\
        P4 & M dwarfs ($>$K7)        & $V < 16$  & $13\,102$ & $4\%$   & primary \\
        P5 & F5--late K dwarfs/subgiants & $V < 13$ & $158\,285$ & $3\%$ & primary \\
        \bottomrule
    \end{tabular}
\end{table}

\paragraph{P2 asteroseismic reference stars.}
P2 is the brightest sub-sample ($V \lesssim 8$), containing $605$ FGK
stars for which \PLATO will deliver the most precise asteroseismic
characterisation.  Their stellar masses and radii will be known to
$\lesssim 1\%$, enabling sub-percent planetary radius measurements.
However, the sample is too small to provide a statistically meaningful
number of habitable-zone planets on both sides of $S_\mathrm{thresh}$,
and it is therefore not analysed in this thesis.

\paragraph{P1 bright FGK dwarfs.}
P1 contains $9\,665$ FGK dwarfs at $V \lesssim 11$, the core science
sample of the \PLATO LOP mission.  These stars will receive
asteroseismology follow-up, and planet radii will be determined to $1\%$.
Although P1 is the highest-priority sample from the mission's perspective,
it still yields too few habitable-zone rocky planets for a robust
demographic test: the geometric transit probability for an Earth analogue
around a G dwarf is only $\approx 0.5\%$, and with $\sim 10^4$ stars the
expected count of detected HZ transiting planets is a few tens at most.
P1 is not analysed further here.

\paragraph{P4 M-dwarf sample.}
P4 contains $13\,102$ M dwarfs spanning $T_\mathrm{eff} = 2926$--$3921\,\mathrm{K}$
and $R_\star = 0.16$--$0.63\,\Rsun$.  M dwarfs are attractive for the
\HZIED test for two reasons.  First, their habitable zones lie at orbital
separations of $0.1$--$0.3\,\mathrm{au}$, giving geometric transit
probabilities of $\sim 5$--$10\%$ — an order of magnitude larger than
for a solar-type star.  Second, the small stellar radii amplify the
transit depth of an Earth-sized planet to several hundred ppm.  The
trade-off is a larger photometric noise per star (P4 stars extend to
$V \approx 15$), giving a per-planet radius uncertainty of $\sigma_R = 4\%$
in the pipeline default.

\paragraph{P5 faint FGK sample.}
P5 is the largest sub-sample, with $158\,285$ FGK stars at $V \lesssim 13$.
It contains $81\,301$ G-type, $65\,358$ F-type, and $11\,585$ K-type
stars.  The reduced photometric precision relative to P1 ($\sigma_R = 3\%$
vs $1\%$) is compensated by the large number of targets: even at a
geometric transit probability of $\lesssim 1\%$, the P5 stellar pool is
expected to yield the largest single-sample catalogue of transiting
rocky habitable-zone planets of any \PLATO sub-sample.

\paragraph{Why P4 and P5.}
P4 and P5 together span the full stellar-mass range accessible to
\PLATO and provide complementary paths to a \HZIED detection.  P4
offers high per-planet transit probability and deep transits from compact
habitable zones; P5 offers statistical power from its large number of
FGK hosts.  Both samples are large enough that the Bioverse generator
can populate them with $\mathcal{O}(100)$ detected habitable-zone planets,
the threshold identified by \textcite{Schlecker2024} for a significant
detection.  The analyses in \Cref{sec:p4_p5_baseline} apply the baseline
runaway greenhouse model to P4, P5, and their combination.


% =======================================================================
\section{From Stellar Catalogue to Detectable Planets}
\label{sec:planet_yield}

Before running the \HZIED hypothesis test, the PIC stellar catalogue must
be converted into a synthetic sample of detected transiting planets.  This
conversion proceeds through three stages: planet generation, detection cuts,
and measurement noise.  The table below traces how the initial pool of
generated planets is reduced at each stage for each \PLATO sample.

\subsection{Planet Generation}
\label{ssec:planet_generation}

Planet populations are drawn from the Bioverse generator using the
\textcite{Bergsten2022} occurrence-rate model.  Each star receives planets
according to the broken power-law radius distribution and period
distribution of that work, with masses assigned via the probabilistic
mass--radius relation of \textcite{ChenKipping2017}.  Stellar ages are
drawn from the individual Gaia/FLAME age estimates attached to each PIC
star (\texttt{age\_mode = `sampled'}), so the age distribution of the
generated population reflects the actual age spread of the \PLATO targets.

\Cref{tab:planet_cuts} shows the cumulative effect of each cut on the
generated populations across all samples.  The generic sample (500 stars,
Gaia-based) is included as the \textcite{Schlecker2024} reference.

\begin{table}[H]
    \centering
    \caption[Planet detection cuts per \PLATO sample]{%
        Cumulative effect of detection cuts on the generated planet
        population for each sample.  Percentages are relative to the
        count in the preceding row (step-wise retention).
        The final yield is the fraction of generated planets that survive
        all cuts and enter the \HZIED analysis.}
    \label{tab:planet_cuts}
    \setlength{\tabcolsep}{5pt}
    \small
    \begin{tabular}{lrrrrr}
        \toprule
        Step & Generic & P1 & P2 & P4 & P5 \\
        \midrule
        $N_\star$ (after quality cuts)
            & 235\,240 & 9\,665 & 605 & 13\,102 & 158\,285 \\
        $N_\mathrm{planets}$ generated
            & 318\,018 & 11\,186 & 696 & 18\,243 & 184\,401 \\
        \quad planets / star
            & 1.35 & 1.16 & 1.15 & 1.39 & 1.16 \\
        \midrule
        After $P \leq 500\,\mathrm{d}$
            & 318\,003 (100\%) & 11\,184 (100\%) & 696 (100\%)
            & 18\,241 (100\%) & 184\,376 (100\%) \\
        After $M \in [0.1,\,2.0]\,\Mearth$
            & 121\,657 (38.3\%) & 3\,222 (28.8\%) & 209 (30.0\%)
            & 6\,540 (35.9\%) & 53\,006 (28.7\%) \\
        After $R > 0.75\,\Rearth$
            & 121\,394 (99.8\%) & 3\,221 (100\%) & 209 (100\%)
            & 6\,529 (99.8\%) & 52\,963 (99.5\%) \\
        Transiting ($|b| \leq 1$)
            & 2\,639 (2.2\%) & 142 (4.4\%) & 8 (3.8\%)
            & 184 (2.8\%) & 2\,217 (4.2\%) \\
        Depth $\geq 80\,\mathrm{ppm}$
            & 2\,535 (96.1\%) & 37 (26.1\%) & 2 (25.0\%)
            & 184 (100\%) & 576 (26.0\%) \\
        $S_\mathrm{abs} \in [10,\,2000]\,\mathrm{W\,m^{-2}}$
            & 1\,541 (60.8\%) & 6 (16.2\%) & 0 (0\%)
            & 82 (44.6\%) & 60 (10.4\%) \\
        \midrule
        Yield (final / generated)
            & 0.485\% & 0.054\% & 0.000\% & 0.449\% & 0.033\% \\
        \bottomrule
    \end{tabular}
\end{table}

\subsection{Detection Cuts and Observable Sample}
\label{ssec:detection_cuts}

The most consequential cut is the transit depth threshold of
$80\,\mathrm{ppm}$, which eliminates Earth-sized planets around
solar-type stars (transit depth $\approx 80\,\mathrm{ppm}$ for
$R_p = 1\,\Rearth$, $R_\star = 1\,\Rsun$).  For the bright FGK
samples P1 and P2, this cut combined with the mass ceiling at
$2.05\,\Mearth$ removes essentially all rocky habitable-zone planets,
leaving $N = 37$ for P1 and $N = 2$ for P2 before the instellation
window cut, and $N = 6$ and $N = 0$ afterwards.  Neither sample
provides enough planets for a rolling-window analysis with $w = 25$
and are therefore excluded from the \HZIED test.

For P4 M~dwarfs, the small stellar radii amplify the transit depth of
an Earth-mass planet to several hundred ppm, so the depth cut retains
all 184 transiting rocky planets (100\%).  The instellation window then
selects those with $S_\mathrm{abs} \in [10,\,2000]\,\mathrm{W\,m^{-2}}$,
yielding the 80--82 planets that enter the analysis.  P5 suffers the
same depth attrition as P1 (26\%), but its large stellar pool produces
58--60 planets in the instellation window.

\Cref{fig:plato_baseline_hzied} shows the observable sample for P4 and
P5 in the radius--instellation plane, colour-coded by classification:
cold-side (blue), hot-side uninflated (grey), and hot-side inflated
(red).  The vertical dashed line marks the injected runaway greenhouse
threshold $S_\mathrm{thresh} = 280\,\mathrm{W\,m^{-2}}$.

\begin{figure}[H]
    \centering
    \includegraphics[width=\textwidth]{plato_baseline_hzied}
    \caption[Detected planet sample for \PLATO P4 and P5 (Turbet baseline)]{%
        Detected transiting rocky planets in the radius--instellation
        plane for the \PLATO P4 (left) and P5 (right) samples.
        Populations generated with the \textcite{Bergsten2022}
        occurrence-rate model; radii inflated with the
        \textcite{Turbet2020} + \textcite{DornLichtenberg2021} model
        at $x_{\mathrm{H_2O}} = 0.005$, $f_{\mathrm{rgh}} = 0.8$;
        stellar ages from Gaia/FLAME (\texttt{age\_mode = `sampled'}).
        Blue: cold-side planets ($S < S_\mathrm{thresh}$).
        Grey: hot-side, uninflated.  Red: hot-side, inflated.
        Error bars show $1\sigma$ measurement uncertainties
        ($\sigma_S = 5\%$; $\sigma_R = 4\%$ P4, $3\%$ P5).
        P4: $N = 80$ ($27$ cold, $8$ hot, $45$ inflated).
        P5: $N = 58$ ($14$ cold, $7$ hot, $37$ inflated).}
    \label{fig:plato_baseline_hzied}
\end{figure}

% Figure produced by:
%   Src/figures/plato_baseline_hzied.py  →  Figures/generated/plato_baseline_hzied.pdf


